% Auto-generated by tables/scripts/make_genrep_corr.py --- do not hand-edit.
\begin{table}[tbp]
\centering
\small
\begin{tabular}{l cc cc}
\toprule
\textbf{Representation quality} & \multicolumn{2}{c}{\textbf{FPSD}} & \multicolumn{2}{c}{\textbf{Designability}} \\
\cmidrule(lr){2-3} \cmidrule(lr){4-5}
 & partial & raw & partial & raw \\
\midrule
Inverse folding (top-1) & +0.35 & +0.13 & +0.79 & +0.58 \\
Dihedral MAE & +0.42 & -0.07 & +0.76 & +0.32 \\
CATH-A & +0.20 & +0.07 & +0.48 & +0.35 \\
\bottomrule
\end{tabular}
\caption{\textbf{Better trunk representations track better generation --- strongly for designability, moderately for FPSD.} Partial Pearson correlation controlling for training step, with the raw value alongside; oriented so a positive value means better representation accompanies better generation. ($n{=}63$ checkpoints pooled over the baseline, the REPA variants, and the random control; generation is a $3$-seed mean, with only the single-seed late-tail checkpoints at $n{=}1$; probe single-seed; $1{,}125$ backbones/seed, $250$ for designability.)}
\label{tab:proteina-genrep-corr}
\end{table}
